\documentclass[../main.tex]{subfiles}
\begin{document}
%==========================================TIÊU ĐỀ TẬP========================================
\begin{table}[h]
    \begin{adjustwidth}{-1cm}{}
        \begin{tabular}{>{\centering\arraybackslash}p{6cm}|>{\raggedright\arraybackslash}p{12.5cm}}
            \multirow{5}{6cm}
            {\\[-40pt]
                \flaregothic\fontsize{20pt}{20pt}\selectfont \centering
                \color{eptype}HISTOIRE\color{black}\\
                \freshbot\fontsize{72pt}{72pt}\selectfont
                \color{epnum}28\\[6pt]
                \cafeteria\fontsize{20pt}{20pt}\selectfont\color{epnum}(D-28)\color{black}
                \\[18pt]
                % \flaregothic\fontsize{18pt}{18pt}\selectfont
                % \color{epattr}TRANSITION
                \color{black}
            }
             &
            {\fotskip\fontsize{18pt}{18pt}\selectfont \ruby{嵐}{あらし}の\ruby{前}{まえ}の\ruby{静}{しず}けさ}
            \\[6pt]
             & {\cafeteria\fontsize{18pt}{18pt}\selectfont The Calm Before the Storm}                                     \\[4pt]
             & {\vnmsans\fontsize{18pt}{18pt}\selectfont Khoảng lặng trước cơn bão}                                       \\[8pt]
             & {\caviar{\fontsize{14pt}{14pt}\selectfont Nhân vật: \emp{kuon2} \emp{kaigou2} \emp{ryuuhou2} \emp{mike1}}}
            \\[8pt]
        \end{tabular}
    \end{adjustwidth}
\end{table}
%=========================================TRUYỆN CHÍNH========================================
\color{black}

\txtbx[2]{
    {\color{vol} \uline{Tóm tắt:}} Sau hai đợt sóng cảm xúc mãnh liệt của chiến dịch Réunion dành cho holoFive và Gen 4, bầu không khí tại ký túc xá Hikawa dần lắng lại trong sự bình yên đầu hè. Thế nhưng, tại văn phòng COVER, những ánh mắt dò hỏi và nỗi niềm khao khát từ các thành viên thế hệ thứ Ba (holoFantasy) bắt đầu dấy lên những làn sóng ngầm. Ý thức rõ vết thương tinh thần sâu sắc và tình trạng bất ổn của Mikeneko trên tầng tám, Kuon cùng hai cô em gái Kaigou và Ryuuhou đã có một cuộc bàn bạc đêm muộn đầy nghiêm túc để vạch ra chiến lược bảo vệ ranh giới an toàn cho ký túc xá. Một đĩa bánh quy còn hơi cháy cạnh cùng mẩu giấy nhắn vụng về nơi khe cửa tầng tám đã tiếp thêm niềm tin cho người Tổng quản lý: dù cơn bão mang tên quá khứ chắc chắn sẽ đến, nhưng sự chữa lành chỉ có thể bắt đầu khi trái tim thực sự sẵn sàng mở lối.
}

\pov{Hikawa Kuon}

\lang{Etsunan}

Đêm cuối tháng Năm.

Cơn mưa rào đầu mùa hạ bất chợt đổ xuống thành phố Kawachi, gột rửa đi lớp bụi bặm ngột ngạt tích tụ suốt những ngày nắng gắt. Từng hạt mưa đập lách tách vào khung cửa kính phòng khách tầng năm, tạo nên một bản hòa tấu trầm buồn nhưng tĩnh mịch đến lạ kỳ.

Tôi ngồi một mình bên bàn trà gỗ, tựa lưng vào thành ghế sofa êm ái, trên tay là tách trà hoa cúc còn bốc khói nghi ngút do Suzuno pha từ chập tối.

Cả căn nhà lúc này đã chìm sâu vào giấc ngủ. Sau những ngày tháng rộn rã tiếng cười đùa của hai chiến dịch Réunion liên tiếp, không gian xung quanh bỗng chốc trở về với vẻ tĩnh lặng nguyên bản vốn có của nó.

Tôi khẽ thở dài, đưa mắt nhìn ra màn mưa mờ ảo ngoài ban công.

Hai cuộc hội ngộ vừa qua thực sự là những phép màu nhiệm màu.

Nhìn Delutaya tìm lại được nụ cười rạng rỡ bên cạnh bốn cô gái NePoLaBo, nhìn Kson gạt bỏ lớp vỏ bọc gai góc để ôm chặt lấy bốn người bạn Gen 4 vào lòng giữa gian kho hàng ngập tràn tiếng cười\dots\ mỗi khoảnh khắc ấy đều như một liều thuốc xoa dịu những vết thương lòng của quá khứ. Các cô gái đã tìm thấy nhau, đã chứng minh cho cả thế giới thấy rằng sợi dây gắn kết chân thành giữa những người đồng đội có thể vượt qua mọi rào cản thời gian và quy tắc nghiệt ngã.

Thế nhưng, chính vì hai cuộc hội ngộ ấy diễn ra quá đỗi trọn vẹn, một áp lực vô hình khác lại bắt đầu đè nặng lên vai tôi.

Tại văn phòng \hll\ những ngày gần đây, bầu không khí đang âm thầm biến chuyển.

Dù các thành viên Gen 4 và Gen 5 đã giữ đúng lời hứa tuyệt đối bảo mật về vị trí của khu ký túc xá nhà Hikawa, nhưng sự thay đổi tích cực đến kinh ngạc trong ánh mắt, nụ cười và năng lượng làm việc của họ là điều không thể nào che giấu nổi. Người ngoài có thể chỉ nghĩ đơn thuần rằng họ vừa vượt qua một nút thắt tâm lý, nhưng với những người gắn bó lâu năm, đó là cả một điều bất thường đáng ngờ.

Đặc biệt là các cô gái của \textit{holoFantasy} - Thế hệ thứ Ba.

Tôi không thể nào quên được ánh mắt của Pekora khi vô tình bắt gặp Kanata và Towa đứng thì thầm to nhỏ với nụ cười rạng rỡ ở góc hành lang phòng thu hồi đầu tuần. Cô thỏ tinh nghịch thường ngày bỗng trở nên trầm ngâm, đôi tai thỏ khẽ cụp xuống với một nét bồn chồn khó tả. Hay như Marine - người thuyền trưởng sở hữu trực giác phụ nữ sắc bén đến đáng sợ - đã hai lần ghé qua bàn làm việc của tôi chỉ để hỏi những câu bâng quơ: ``Kuon-senpai dạo này có vẻ hay đi về Kawachi muộn nhỉ?'', rồi lại lén nhìn tôi với ánh mắt chan chứa sự hoài nghi lẫn hy vọng.

Họ đang tìm kiếm một câu trả lời.

Họ vẫn luôn mang trong tim một khoảng trống mênh mông kể từ biến cố chấn động vào cuối tháng Hai năm 2022.

Cái ngày mà Uruha Rushia chính thức rời khỏi công ty.

Một sự ra đi không có buổi lễ tốt nghiệp êm đẹp như Coco, không có những lời tạm biệt trong nước mắt hòa cùng tiếng vỗ tay như Aloe. Đó là một cơn cuồng phong tàn khốc của dư luận, những bản hợp đồng bị xé rách, những cáo buộc, những lời thóa mạ cay độc nhất từ không gian mạng dội thẳng vào tâm can một cô gái nhỏ bé mang tâm hồn mong manh như thủy tinh.

Và kết quả của cơn bão ấy\dots\ hiện đang trú ngụ ngay trên tầng tám của tòa nhà này.

\uline{Mikeneko.}

\dots

``Anh lại ngồi trầm ngâm một mình nữa rồi đấy, anh Kuon.''

Một giọng nói điềm tĩnh, quen thuộc vang lên phá vỡ dòng suy nghĩ miên man của tôi.

Tôi ngẩng đầu lên. Kaigou đang bước vào phòng khách, trên người vẫn mặc chiếc áo thun xám đơn giản với chân váy đen. Mái tóc đen được cột bằng buộc tóc cầu vồng còn hơi ẩm ướt sau khi tắm, quanh cổ còn đang quàng một chiếc khăn tắm. Em gái song sinh của tôi kéo chiếc ghế đối diện ngồi xuống, đặt lên bàn một chiếc đĩa đựng vài lát táo đã gọt sẵn.

``Sao chưa ngủ thế?'' tôi khẽ hỏi, nhấp một ngụm trà. ``Mai em còn phải lên công ty sớm để hỗ trợ bên \hls\ cơ mà?''

``Anh tưởng em ngủ được khi thấy anh ngồi thở ngắn than dài từ lúc mười một giờ đêm đến giờ chắc?'' Kaigou khoanh hai tay trước ngực, đôi mắt xanh sắc sảo nhìn xoáy vào tôi. ``Lại đang nghĩ về chuyện của Thế hệ thứ Ba đúng không?''

Tôi im lặng một thoáng, rồi khẽ gật đầu: ``Ừ. Cũng chẳng thể nào giấu nổi em nhỉ.''

``Chuyện đó là tất yếu thôi,'' Kaigou thở dài, giọng nói chùng xuống mang vẻ nghiêm túc hiếm thấy. ``Chiến dịch Réunion của Gen 5 và Gen 4 thành công, người ngoài nhìn vào thì vui đấy, nhưng ranh giới bảo mật của chúng ta đang mỏng manh hơn bao giờ hết. Văn phòng bên kia bắt đầu râm ran tin đồn rồi. Sáng nay lúc ở căng-tin, Noel-san đã vô tình hỏi em xem ký túc xá nhà mình dạo này có ai chuyển đến không.''

Tôi giật mình, đặt chén trà xuống: ``Noel hỏi thẳng em luôn à?''

``Vâng. Em ấy hỏi với vẻ mặt ngập ngừng lắm, trông tội nghiệp vô cùng. Em chỉ đành giả ngơ bảo là nhà vẫn chỉ có mấy anh em mình và bọn trẻ thôi. Nhưng ánh mắt của Noel-san\dots\ em biết em ấy không hoàn toàn tin đâu. Trực giác của những người cùng chung một thế hệ nhạy bén lắm, nhất là khi họ đã cùng nhau đi qua những năm tháng rực rỡ nhất.''

``Đúng vậy\dots'' tôi day day sống mũi, cảm thấy thái dương bắt đầu nhức nhối. ``Nhưng Mikeneko thì không giống như Aloe hay Coco. Em ấy chưa hề sẵn sàng.''

``Chuẩn không cần chỉnh luôn đấy Onii!''

Một giọng nói lanh lảnh, mang đậm vẻ bộc trực bất ngờ vang lên từ phía cửa bếp.

Ryuuhou bước ra, tay cầm một hộp sữa chua dâu và chiếc thìa nhựa. Em ấy vẫn khoác chiếc áo khoác nỉ ngoại cỡ của ban nhạc underground, đầu đội chiếc mũ lưỡi trai ngược tinh nghịch.

``Ryuuhou? Sao em còn thức?'' Kaigou nheo mắt lườm cô em út. ``Mười hai giờ đêm rồi đấy cô nương! Mai không phải đi học à?''

``Mai thứ Bảy mà chị Kaigou, trường em được nghỉ!'' Ryuuhou bĩu môi đáp trả, rồi thoăn thoắt đi tới ngồi phịch xuống tấm thảm lông ngay cạnh đầu gối tôi. Con bé múc một muỗng sữa chua cho vào miệng, nhai nhóp nhép rồi hướng đôi mắt sáng ngời về phía hai chúng tôi: ``Em vừa nghe hết câu chuyện rồi nha. Về chị Mikeneko trên tầng tám đúng không?''

Tôi nhìn cô em gái tomboy của mình, khẽ cười trừ: ``Cái tai thính của dân chơi nhạc underground có khác. Em nghĩ sao về chuyện này?''

Ryuuhou buông chiếc muỗng xuống, nét mặt bỗng chốc trở nên nghiêm túc khác hẳn ngày thường. Con bé dựa lưng vào ghế sofa, khoanh hai tay lại, cất giọng rành rọt: ``Em nói thật lòng nhé anh Kuon, nếu bây giờ mà để các chị bên Gen 3 đùng đùng kéo đến đây đòi gặp, thì đó không phải là một cuộc hội ngộ đâu. Đó sẽ là một vụ `thảm sát tâm lý' đấy.''

Lời nhận xét thẳng thừng như một gáo nước lạnh tạt thẳng vào căn phòng, khiến cả tôi và Kaigou đều sững sờ im lặng.

``Em lăn lộn trong mấy ban nhạc đường phố từ năm mười hai tuổi rồi,'' Ryuuhou tiếp tục, ánh mắt nhìn thẳng vào tôi không hề né tránh. ``Em từng thấy không biết bao nhiêu người bị hủy hoại chỉ vì một bài phốt trên mạng hay sự phản bội của bạn bè. Vết thương do đòn roi thì vài tuần là liền sẹo, nhưng vết thương do bị hàng chục ngàn người cùng lúc quay lưng, chửi rủa, gọi mình là kẻ dối trá\dots\ nó là thứ độc dược ăn mòn tâm trí từng ngày. Chị Mikeneko mang đúng thứ vết thương đó.''

Con bé thở dài, giọng nói chùng xuống đầy xót xa: ``Anh với chị Kaigou ban ngày đi làm suốt, chỉ có em, chị Suzuno với bé Koneko là hay ở nhà nhiều nhất thôi. Có những hôm trời mưa to sấm chớp thế này, em lén đi ngang qua tầng tám, vẫn nghe thấy tiếng chị ấy thở gấp rồi khóc thút thít trong phòng. Mỗi lần nghe thấy tiếng chuông cửa hay tiếng bước chân người lạ ngoài cầu thang, chị ấy đều giật thót mình, rụt cổ lại như một con mèo bị dính bẫy vậy. Chị ấy sợ thế giới bên ngoài, sợ cả những người từng thân thiết vì trong tiềm thức, chị ấy luôn nghĩ mình là kẻ tội đồ đã làm ô uế cái tên của nhóm.''

Kaigou khẽ cắn môi, những ngón tay siết chặt lại trên mặt bàn: ``Ryuuhou nói đúng đấy anh Kuon. Aloe ngày trước rời đi vì áp lực quá tải, nhưng trong lòng em ấy vẫn luôn khao khát được hát, được gặp lại NePoLaBo. Coco thì bản lĩnh ngút trời, ra đi trong thế ngẩng cao đầu để bảo vệ đàn em. Còn Mikeneko\dots\ niềm tin của con bé vào con người đã sụp đổ hoàn toàn vào cái ngày bản hợp đồng ấy bị chấm dứt. Con bé chỉ vừa mới bắt đầu học cách hít thở lại bình thường trong căn nhà này thôi.''

Tôi lặng người lắng nghe từng lời của hai cô em gái. Trái tim tôi thắt lại khi nghĩ đến hình bóng gầy gò, cô độc của Mikeneko.

Đúng như lời Ryuuhou nói. Những bước tiến gần đây của Mikeneko - từ việc lặng lẽ để lại cốc nước ấm đầu giường khi tôi bị sốt, cho đến việc rụt rè ngồi ở góc phòng khách xem trọn bộ phim hoạt hình gia đình trong đêm chiếu phim - tất cả đều là những tia hy vọng vô cùng mong manh. Đó là những chồi non vừa mới nhú lên từ đống tro tàn đổ nát. Nếu lúc này chúng ta vội vã đưa một cơn gió mạnh mang tên ``holoFantasy'' ập tới, những chồi non ấy sẽ gãy vụn ngay lập tức mà không bao giờ có thể đâm chồi được nữa.

``Vậy nên,'' Ryuuhou gõ nhẹ đầu ngón tay lên mặt bàn, đưa ra kết luận dứt khoát: ``Tạm thời tuyệt đối không được để các chị Gen 3 biết chỗ này. Cho dù chị Marine có tinh quái thế nào, chị Pekora có bồn chồn ra sao, anh Kuon cũng phải làm tròn vai một bức tường sắt kiên cố ở công ty. Không được hé răng nửa lời!''

Kaigou gật đầu tán thành: ``Em đồng ý với Ryuu. Chúng ta phải thiết lập một vùng đệm an toàn tuyệt đối cho tầng tám. Chỉ khi nào Mikeneko tự tay mở toang cánh cửa phòng mình, tự mình nói rằng con bé muốn gặp lại những người đồng đội cũ, thì lúc đó chiến dịch Réunion mới được phép bắt đầu. Bằng không, bất kỳ sự can thiệp sớm nào cũng sẽ là một tội ác.''

Nhìn hai cô em gái - một người sắc sảo, lý trí với tư duy logic của một con người ở thế giới cũ, một người từng trải, thấu hiểu lòng người bằng bản năng nhạy bén của giới âm nhạc đường phố - tôi cảm thấy trong lòng dâng lên một niềm cảm kích và tự hào khôn xiết.

Gia đình Hikawa của tôi thực sự đã lớn khôn. Họ không chỉ là những đứa em gái cần tôi che chở, mà giờ đây đã trở thành những điểm tựa vững chãi, cùng tôi gánh vác những trọng trách nặng nề nhất.

Tôi mỉm cười, đưa tay xoa nhẹ đầu Ryuuhou khiến con bé giật mình phụng phịu né ra, rồi nhìn sang Kaigou: ``Anh hiểu rồi. Cảm ơn hai đứa. Anh hứa sẽ giữ vững phòng tuyến ở văn phòng. Cho dù có phải đóng vai một ông Tổng quản lý lạnh lùng, vô cảm trước mặt Thế hệ thứ Ba, anh cũng sẽ không để bất kỳ ai làm xáo trộn cuộc sống bình yên của Mikeneko lúc này.''

``Thế mới là anh trai của bọn em chứ!'' Ryuuhou nhe răng cười toe toét, vỗ ngực đôm đốp.

``Thôi, muộn lắm rồi đấy,'' Kaigou đứng dậy, giật lấy hộp sữa chua trên tay Ryuuhou. ``Đi đánh răng rồi đi ngủ ngay cho tôi cô nương! Sáng mai mà ngủ nướng đến trưa là cắt tiền tiêu vặt đấy nhé!''

``Ớ! Chị Kaigou độc tài! Em còn một miếng cuối mà!'' Ryuuhou làu bàu giãy nảy lên, nhưng rồi cũng ngoan ngoãn lon ton chạy theo chị gái về phía hành lang phòng ngủ.

Không gian phòng khách lại trở về với sự tĩnh lặng ban đầu.

\ct

Một giờ sáng.

Mưa ngoài trời đã bắt đầu ngớt dần, chỉ còn lại những giọt nước đọng trên mái hiên tí tách rơi xuống nền xi-măng.

Tôi tắt đèn phòng khách, nhẹ nhàng rảo bước về phía cầu thang bộ. Thay vì đi về phòng ngủ của mình ở tầng năm, đôi chân tôi lại vô thức dẫn lối, từng bước chầm chậm leo lên những bậc thang dẫn lên tầng cao nhất của tòa nhà.

Tầng tám.

Hành lang tầng tám luôn chìm trong ánh đèn vàng dịu nhẹ, mờ ảo và tĩnh mịch hơn bất kỳ nơi nào trong khu ký túc xá. Nơi đây vốn được thiết kế biệt lập hoàn toàn, dành riêng cho hai chị em Mikeneko và Koneko để đảm bảo sự riêng tư tuyệt đối cho quá trình hồi phục tâm lý.

Tôi bước nhẹ trên nền gỗ, cố gắng không phát ra bất kỳ tiếng động nào.

Căn phòng ở cuối hành lang vẫn sáng đèn. Qua khe hở mỏng manh dưới chân cánh cửa gỗ sồi sẫm màu, ánh sáng xanh nhạt từ màn hình máy tính vẫn le lói hắt ra ngoài.

Em ấy vẫn chưa ngủ.

Tôi dừng chân cách cánh cửa chừng hai bước. Bàn tay tôi bất giác nâng lên giữa không trung, những ngón tay khẽ co lại, định gõ nhẹ ba tiếng quen thuộc để hỏi thăm em như thường lệ.

Thế nhưng, bàn tay tôi bỗng khựng lại.

Lời nói của Ryuuhou và Kaigou lúc nãy lại vang vọng trong tâm trí tôi: ``Chỉ khi nào con bé tự tay mở toang cánh cửa phòng mình\dots''

Sự quan tâm quá mức đôi khi cũng là một thứ áp lực vô hình. Một người quản lý tốt không phải là người luôn luôn xông vào thế giới của người khác để ban phát sự giúp đỡ, mà là người biết kiên nhẫn đứng đợi ở bên ngoài, làm một ngọn hải đăng vững chãi để khi đối phương sẵn sàng bước ra, họ biết chắc chắn rằng mình luôn có một bến đỗ an toàn.

Tôi thở dài một hơi thật khẽ, từ từ hạ tay xuống.

``\textit{Chưa phải lúc\dots}''

Tôi khẽ xoay người, chuẩn bị bước xuống cầu thang thì bất chợt\dots

*CẠCH*

Một tiếng động kim loại cực kỳ nhỏ vang lên từ chốt khóa cửa.

Tôi giật mình khựng lại, quay đầu nhìn lại. Cánh cửa phòng vốn luôn đóng chặt kín mít suốt bao ngày qua\dots\ lúc này đang khẽ hé mở một khe nhỏ chừng vài centimet.

Một luồng gió mát mang theo mùi hương vani nhè nhẹ từ trong phòng thoảng ra ngoài.

Không có ai bước ra. Tôi chỉ thoáng thấy một bóng dáng nhỏ nhắn trong bộ đồ ngủ màu đen nhanh như cắt rụt lại vào sau cánh cửa, hệt như một chú mèo con sợ người lạ đang lén lút quan sát qua khe hở.

Tôi đứng yên bất động, không dám tiến lại gần vì sợ làm em ấy giật mình hoảng sợ.

Và rồi, từ dưới khe cửa hẹp ấy, một bàn tay nhỏ nhắn, trắng bợt khẽ thò ra ngoài. Đôi bàn tay ấy cẩn thận đẩy nhẹ một chiếc khay nhựa màu trắng ra giữa sàn hành lang, rồi vội vã rụt phắt lại vào trong.

*CẠCH*

Cánh cửa lại khép hờ lại, chốt cửa không hề khóa.

Tôi chớp chớp mắt kinh ngạc. Tim tôi đập rộn lên từng nhịp thổn thức. Tôi bước lại gần chiếc khay nhựa, cúi người xuống nhìn kỹ.

Trên chiếc khay là một chiếc đĩa sứ nhỏ đựng bốn chiếc bánh quy bơ nướng. Những chiếc bánh có hình dáng tròn méo vụng về, vài góc cạnh hơi bị cháy sém màu nâu sẫm - rõ ràng là thành quả của một buổi tập làm bánh đầy bỡ ngỡ giữa Mikeneko và cô em gái Koneko vào buổi chiều.

Và đặt ngay bên cạnh chiếc đĩa là một mẩu giấy nhớ màu hồng phấn cắt hình chiếc đầu mèo.

Trên mẩu giấy, ai đó đã dùng bút dạ vẽ một hình mặt mèo chibi đang mỉm cười ngượng nghịu với hai má ửng hồng. Phía dưới hình vẽ là một dòng chữ viết tay nắn nót nhưng những nét mực lại hơi run rẩy, như thể người viết đã phải gom hết can đảm của mình mới có thể đặt bút:

\begin{center}
    \textit{``Cảm ơn anh Kuon\dots\ vì đã luôn bảo vệ nơi này. Bánh hơi cháy một chút\dots\ nhưng anh ăn tạm nhé.''}
\end{center}

Tôi nhìn chằm chằm vào dòng chữ ấy trong vài giây.

Một cảm giác ấm áp đến nghẹn ngào bất chợt dâng trào từ tận đáy lòng, lan tỏa khắp lồng ngực tôi. Sống mũi tôi cay cay, khóe môi không kìm được mà cong lên một nụ cười nhẹ nhõm, hạnh phúc nhất trần đời.

Em ấy không hề trốn chạy.

Em ấy biết hết tất cả. Nó biết những nỗ lực thầm lặng của tôi, của Kaigou, của Ryuuhou và cả gia đình Hikawa suốt bao ngày qua để giữ cho hai chị em họ một khoảng trời bình yên. Chiếc bánh quy vụng về này, mẩu giấy nhắn run rẩy này\dots\ chính là cánh cửa đầu tiên mà Mikeneko đã tự tay mở ra với thế giới bên ngoài.

Tôi quỳ một bên gối xuống sàn, nhẹ nhàng nhấc chiếc khay lên tay.

Tôi cầm một chiếc bánh quy đưa lên miệng cắn một miếng nhỏ. Vị ngọt bùi của bơ hòa quyện với vị đăng đắng nhẹ của phần rìa bị cháy lan tỏa trên đầu lưỡi - đối với tôi lúc này, đó là món bánh ngon nhất mà tôi từng được thưởng thức trong suốt hai mươi lăm năm cuộc đời.

Tôi ngước mắt nhìn lên khe cửa vẫn đang khép hờ, cất giọng thật trầm ấm, dịu dàng, đủ để người bên trong có thể nghe thấy rõ từng lời: ``Bánh ngon lắm Mike-chan. Ngọt và thơm lắm. Cảm ơn hai chị em nhé.''

Tôi dừng lại một chút, rồi khẽ nói thêm: ``Mưa tạnh rồi đấy. Mai trời sẽ nắng đẹp trở lại thôi. Hai đứa ngủ sớm đi nhé.''

Phía sau cánh cửa gỗ, tôi nghe thấy một tiếng sột soạt rất khẽ, tựa như tiếng một chú mèo con vừa cuộn tròn mình lại trong chiếc chăn ấm, mang theo một tiếng thở phào nhẹ nhõm vô cùng bình yên.

Ánh đèn hắt qua khe cửa khẽ chớp tắt một nhịp, như một lời chúc ngủ ngon không thành tiếng.

\ct

Tôi cầm chiếc khay bước lên sân thượng tầng thượng của tòa nhà.

Mưa đã tạnh hẳn. Những đám mây đen kịt đã bị làn gió đêm quét sạch về phía chân trời, để lộ ra bầu trời đêm Kawachi trong vắt, lấp lánh muôn ngàn vì tinh tú rạng ngời. Không khí sau cơn mưa trong lành và mát rượi, thoang thoảng mùi ngai ngái nồng nàn của đất mẹ và hương hoa dạ lý hương ven đường.

Tôi tựa hai tay vào lan can sắt, phóng tầm mắt nhìn về phía xa xăm.

Phía chân trời đằng đông, ánh đèn rực rỡ của thủ đô Tokyo và thành phố Điện tử vẫn đang lấp lánh như một dải ngân hà rực rỡ dưới mặt đất. Ở nơi đó, cơn sóng ngầm mang tên \textit{holoFantasy} chắc chắn vẫn đang âm ỉ cuộn trào. Sẽ đến một ngày, Pekora, Marine, Flare và Noel sẽ tìm ra manh mối, sẽ tìm đến trước cánh cổng của ngôi nhà này để đòi lại câu trả lời cho những nỗi đau bị chôn giấu suốt bao tháng ngày qua.

Đó sẽ là một cơn bão lớn nhất, khốc liệt nhất mà tôi từng phải đối mặt kể từ ngày tiếp quản vai trò Tổng quản lý.

Thế nhưng, nhìn những vì sao lấp lánh trên cao, tôi không còn cảm thấy hoang mang hay lo sợ nữa.

Bên cạnh tôi có một gia đình tuyệt vời - có một Kaigou bản lĩnh, một Ryuuhou thấu hiểu lòng người, một Suzuno ấm áp và một Naomi thuần khiết. Và ngay dưới chân tôi, trong căn phòng nhỏ ở tầng tám, hạt mầm của sự hàn gắn đã bắt đầu cựa mình thức giấc.

Mọi vết thương rồi sẽ lành, miễn là chúng ta cho nó đủ thời gian và sự kiên nhẫn.

Tôi cầm chiếc bánh quy thứ hai lên ăn, mỉm cười ngắm nhìn bầu trời đêm lộng gió.

Chưa phải lúc\dots\ nhưng khi cơn bão ấy thực sự ập đến, chúng tôi nhất định sẽ sẵn sàng.

\end{document}